The data are insufficient. Put $A=(0,0),B=(w,0),D=(0,6),P=(\dfrac{w}{3},0),Q=(x,y)$. The condition $AQ=6$ gives $x^2+y^2=36$. The upper circle locus for $\angle PQB=45^\circ$ has equation $x^2+y^2-\dfrac{4w}{3}x-\dfrac{2w}{3}y+\dfrac{w^2}{3}=0$, hence $w^2-2w(2x+y)+108=0$. One valid example is $w=18,Q=(\dfrac{18}{5},\dfrac{24}{5})$, area 108. Another is $w=12,Q=(\dfrac{42-3\sqrt{31}}{10},\dfrac{21+6\sqrt{31}}{10})$, area 72. Both points are strictly inside their rectangles and satisfy the equations and the acute 45-degree angle. Thus the area is not uniquely determined.
\begin{asy}
import graph;import geometry;usepackage("amsmath");size(7cm);
pair A=(0,0),B=(18,0),C=(18,6),D=(0,6),P=(6,0),Q=(3.6,4.8);draw(A--B--C--D--cycle);draw(A--Q^^P--Q--B,blue);label("$A$",A,(-1,-1));label("$B$",B,(1,-1));label("$C$",C,(1,1));label("$D$",D,(-1,1));label("$P$",P,(0,-1));label("$Q$",Q,(0,1));label("$6$",(A+Q)/2,(-1,0));label("example: area $108$",(10,5));
\end{asy}

